% Supported v0.5 artist-statement example.
\documentclass[type=artist]{careerdossier-statement}
% examples/profiles/profile-academic.tex
%
% Shared profile metadata used by the academic-CV example and the corresponding
% industry-resume smoke fixture. Keeping it independent of document layout
% proves that both document families consume the same \CDossierSetup interface.

\CDossierSetup{
  name     = {Ada Lovelace},
  headline = {Researcher in Analytical Computing},
  email    = {ada.lovelace@example.com},
  location = {Toronto, Ontario, Canada},
  website  = {example.com/ada-lovelace},
  %% The bare Google Scholar identifier, expanded to
  %% scholar.google.com/citations?user=ada-example for both the displayed text
  %% and the link. Pasting the whole address works too and is left as written.
  %% See docs/API.md, "Accepted forms for the web-profile keys".
  scholar  = {ada-example},
  orcid    = {0000-0002-1825-0097},
  affiliation = {Example University},
}

\CDossierStatementSetup{
  subtitle={Material traces of computation and memory},
  application-context={Application to the PhD in Visual Arts, Example University}
}
\begin{document}
\MakeCDossierStatementHeader

\section*{Artistic Practice}

I work across print, installation, and small computational systems to examine
how acts of measurement become stories about a place. My practice begins with
records that appear objective---sensor readings, maps, maintenance logs, image
metadata---and asks what they preserve, what they omit, and who is able to
interpret them. I translate these records into material situations that slow
down the encounter with data and return attention to uncertainty, labour, and
memory.

The work is not an illustration of technology. I am interested in the distance
between a system's promise of precision and the partial conditions under which
it operates. A reading changes with weather, placement, calibration, and care;
a map changes with the categories selected before a line is drawn. By making
those conditions visible, I invite viewers to consider measurement as a social
and material practice rather than a neutral view from nowhere.

\section*{Methods and Materials}

My process combines field observation, archival research, code, and repetitive
hand work. I collect modest datasets from walks and conversations, then build
scripts that reorganize them by interval, absence, or contradiction rather than
by a single summary value. The resulting patterns guide etched plates, layered
prints, projected text, or arrangements of found material. Moving between code
and hand processes lets error remain perceptible: a registration shift, a gap
in a record, or an unexpected density can redirect the work.

In the installation \emph{Weather Index}, a series of paper strips carries
hourly environmental readings gathered around a building scheduled for
renovation. The values determine the spacing of marks but not their colour,
which is mixed from dust collected at the site. Visitors can follow the numeric
sequence, yet the fragile surface and uneven pigment resist the appearance of a
complete archive. The piece treats deterioration not as noise to remove but as
evidence of the building's changing use.

Another project, \emph{Unresolved Coordinates}, begins with public maps whose
categories changed over several decades. I print incompatible boundaries on
translucent sheets and suspend them so that alignment depends on the viewer's
position. Short fragments of oral history are available as text and audio near
the work, not as labels that settle the image but as accounts that complicate
the administrative record.

\section*{Context and Development}

My practice is informed by artists who use archives critically, by feminist
approaches to situated knowledge, and by community-based methods that treat
interpretation as a relationship rather than extraction. I am cautious about
turning another person's experience into raw material. Projects therefore begin
with conversations about purpose, access, attribution, and what should not be
collected. When work involves community knowledge, participants can review how
their contributions are represented and determine whether material remains
public after an exhibition.

Recent work has moved from static records toward systems that change in the
presence of viewers. I am developing low-energy devices that reveal their own
sampling and failure states rather than hiding them behind a seamless interface.
The challenge is to keep the technical mechanism legible without allowing it to
dominate the sensory experience. I want the viewer to notice both the event and
the conditions that made it visible.

\section*{Research Direction and Programme Fit}

At Example University, I hope to deepen this inquiry through sustained studio
research on material computation, critical mapping, and participatory archives.
Access to print facilities would allow larger experiments with layered and
time-based surfaces, while dialogue across visual arts, geography, and digital
humanities would strengthen the ethical and historical grounding of the work.

The doctoral project would produce a connected body of installations and a
written account of how transparent computational processes can function as
artistic form. I aim to develop work that neither celebrates nor rejects data,
but makes room to experience its incompleteness---and to ask what kinds of
attention, responsibility, and collective memory might grow from that space.

\end{document}
